% GENERATED FILE. Edit canonical lessons, then run npm run generate:books.
% canonical-source-hash: fnv1a64:6ea7e09f02779f05
% canonical-lessons: KA-C270-bhugola, KA-C270-pada, KA-C270-vakya, KA-C270-ksaurika, KA-C270-dhobi

\chapter{Geography, Word, Sentence, Barber, Washerman}
\label{ch:ka-geography-word-sentence-barber-washerman}
\hlchaptermodality{270}

\begin{chapteropening}
\textbf{By the end of this chapter:} \emph{I can say geography, a word, a sentence, a barber and a washerman.}

Complete the last lesson of chapter 270: I can say geography, a word, a sentence, a barber and a washerman.
\end{chapteropening}

\section[\texorpdfstring{\kn{ಭೂಗೋಳ} (bhūgōḷa)}{ಭೂಗೋಳ (bhūgōḷa)}]{\kn{ಭೂಗೋಳ} (bhūgōḷa) — geography}

\label{lesson:KA-C270-bhugola}



\begin{quote}
Before the new one: say the Kannada for information, then the Kannada for a mark.
\end{quote}

\subsection*{You'll want to know: \kn{ಭೂಗೋಳ}}
\textbf{\kn{ಭೂಗೋಳ}} — \emph{bhūgōḷa} — \textquotedblleft{}geography\textquotedblright{}.

\begin{grammarlens}[title={What you've built}]
One more word for this chapter.
\end{grammarlens}

\subsection*{Your turn}
\begin{itemize}
\raggedright
  \item \emph{Say it:} \emph{bhūgōḷa}
  \item \emph{Say it:} \emph{bhūgōḷa}, once more
  \item \emph{Say it:} say \emph{bhūgōḷa}
  \item \emph{Recall:} say the Kannada for information, then the Kannada for a mark, then say \emph{bhūgōḷa} again
\end{itemize}

\begin{tcolorbox}[breakable,colback=teal!4,colframe=teal!35!black,title={Before you move on}]
What does \kn{ಭೂಗೋಳ} mean? (\textquotedblleft{}Geography\textquotedblright{}.) Say it once more.
\end{tcolorbox}
\section[\texorpdfstring{\kn{ಪದ} (pada)}{ಪದ (pada)}]{\kn{ಪದ} (pada) — a word}

\label{lesson:KA-C270-pada}



\begin{quote}
Before the new one: say the Kannada for a mark, then the Kannada for geography.
\end{quote}

\subsection*{You'll want to know: \kn{ಪದ}}
\textbf{\kn{ಪದ}} — \emph{pada} — \textquotedblleft{}a word\textquotedblright{}.

\begin{grammarlens}[title={What you've built}]
One more word for this chapter.
\end{grammarlens}

\subsection*{Your turn}
\begin{itemize}
\raggedright
  \item \emph{Say it:} \emph{pada}
  \item \emph{Say it:} \emph{pada}, once more
  \item \emph{Say it:} say \emph{pada}
  \item \emph{Recall:} say the Kannada for a mark, then the Kannada for geography, then say \emph{pada} again
\end{itemize}

\begin{tcolorbox}[breakable,colback=teal!4,colframe=teal!35!black,title={Before you move on}]
What does \kn{ಪದ} mean? (\textquotedblleft{}A word\textquotedblright{}.) Say it once more.
\end{tcolorbox}
\section[\texorpdfstring{\kn{ವಾಕ್ಯ} (vākya)}{ವಾಕ್ಯ (vākya)}]{\kn{ವಾಕ್ಯ} (vākya) — a sentence}

\label{lesson:KA-C270-vakya}



\begin{quote}
Before the new one: say the Kannada for geography, then the Kannada for a word.
\end{quote}

\subsection*{You'll want to know: \kn{ವಾಕ್ಯ}}
\textbf{\kn{ವಾಕ್ಯ}} — \emph{vākya} — \textquotedblleft{}a sentence\textquotedblright{}.

\begin{grammarlens}[title={What you've built}]
One more word for this chapter.
\end{grammarlens}

\subsection*{Your turn}
\begin{itemize}
\raggedright
  \item \emph{Say it:} \emph{vākya}
  \item \emph{Say it:} \emph{vākya}, once more
  \item \emph{Say it:} say \emph{vākya}
  \item \emph{Recall:} say the Kannada for geography, then the Kannada for a word, then say \emph{vākya} again
\end{itemize}

\begin{tcolorbox}[breakable,colback=teal!4,colframe=teal!35!black,title={Before you move on}]
What does \kn{ವಾಕ್ಯ} mean? (\textquotedblleft{}A sentence\textquotedblright{}.) Say it once more.
\end{tcolorbox}
\section[\texorpdfstring{\kn{ಕ್ಷೌರಿಕ} (kṣaurika)}{ಕ್ಷೌರಿಕ (kṣaurika)}]{\kn{ಕ್ಷೌರಿಕ} (kṣaurika) — a barber}

\label{lesson:KA-C270-ksaurika}



\begin{quote}
Before the new one: say the Kannada for a word, then the Kannada for a sentence.
\end{quote}

\subsection*{You'll want to know: \kn{ಕ್ಷೌರಿಕ}}
\textbf{\kn{ಕ್ಷೌರಿಕ}} — \emph{kṣaurika} — \textquotedblleft{}a barber\textquotedblright{}.

\begin{grammarlens}[title={What you've built}]
One more word for this chapter.
\end{grammarlens}

\subsection*{Your turn}
\begin{itemize}
\raggedright
  \item \emph{Say it:} \emph{kṣaurika}
  \item \emph{Say it:} \emph{kṣaurika}, once more
  \item \emph{Say it:} say \emph{kṣaurika}
  \item \emph{Recall:} say the Kannada for a word, then the Kannada for a sentence, then say \emph{kṣaurika} again
\end{itemize}

\begin{tcolorbox}[breakable,colback=teal!4,colframe=teal!35!black,title={Before you move on}]
What does \kn{ಕ್ಷೌರಿಕ} mean? (\textquotedblleft{}A barber\textquotedblright{}.) Say it once more.
\end{tcolorbox}
\section[\texorpdfstring{\kn{ಧೋಬಿ} (dhōbi)}{ಧೋಬಿ (dhōbi)}]{\kn{ಧೋಬಿ} (dhōbi) — a washerman, a laundryman}

\label{lesson:KA-C270-dhobi}



\begin{quote}
Before the new one: say the Kannada for a sentence, then the Kannada for a barber.
\end{quote}

\subsection*{You'll want to know: \kn{ಧೋಬಿ}}
\textbf{\kn{ಧೋಬಿ}} — \emph{dhōbi} — \textquotedblleft{}a washerman, a laundryman\textquotedblright{}.

\begin{grammarlens}[title={What you've built}]
One more word for this chapter.
\end{grammarlens}

\subsection*{Your turn}
\begin{itemize}
\raggedright
  \item \emph{Say it:} \emph{dhōbi}
  \item \emph{Say it:} \emph{dhōbi}, once more
  \item \emph{Say it:} say \emph{dhōbi}
  \item \emph{Recall:} say the Kannada for a sentence, then the Kannada for a barber, then say \emph{dhōbi} again
\end{itemize}

\begin{tcolorbox}[breakable,colback=teal!4,colframe=teal!35!black,title={Before you move on}]
What does \kn{ಧೋಬಿ} mean? (\textquotedblleft{}A washerman, a laundryman\textquotedblright{}.) Say it once more.
\end{tcolorbox}
